% ===== PRÉAMBULE CANONIQUE — à coller VERBATIM en tête de CHAQUE .tex =====
\documentclass[11pt,a4paper]{article}
\usepackage[utf8]{inputenc}\usepackage[T1]{fontenc}\usepackage{lmodern}
\usepackage[french]{babel}
\usepackage{amsmath,amssymb,mathtools}\usepackage{icomma}
\usepackage[version=4]{mhchem}          % \ce{...} : équations & formules
\usepackage{siunitx}\sisetup{locale=FR} % \qty{}{}, \si{}, \ang{}
\usepackage{chemfig}                    % \chemfig{...} : structures développées
\usepackage{xcolor}\usepackage{colortbl}
\usepackage{tikz}\usepackage{pgfplots}\pgfplotsset{compat=1.18}
\usepackage[most]{tcolorbox}\usepackage{enumitem}\usepackage{array,booktabs}
\usepackage[a4paper,margin=2.2cm]{geometry}
\usetikzlibrary{arrows.meta,calc,angles,quotes,positioning,patterns,backgrounds,shapes.geometric}
\definecolor{primary}{RGB}{222,49,99}\definecolor{primarydark}{RGB}{150,22,55}
\definecolor{defcol}{RGB}{0,114,178}\definecolor{methcol}{RGB}{52,92,148}
\definecolor{excol}{RGB}{0,135,90}\definecolor{attncol}{RGB}{200,40,55}
\tcbset{boxbase/.style={enhanced,breakable,boxrule=0.9pt,arc=2.5pt,
  left=3mm,right=3mm,top=2.3mm,bottom=2.3mm,fonttitle=\bfseries,coltitle=white}}
% --- boîtes TUTORIEL (noms EXACTS, ne pas renommer) ---
\newtcolorbox{cle}[1][]{boxbase,colback=defcol!6,colframe=defcol,
  title={\textsc{À retenir}\ifx\relax#1\relax\relax\else\textmd{ : \itshape #1}\fi}}
\newtcolorbox{methode}[1][]{boxbase,colback=methcol!7,colframe=methcol,sharp corners,
  title={\textsc{Méthode}\ifx\relax#1\relax\relax\else\textmd{ --- \itshape #1}\fi}}
\newtcolorbox{exemplebox}[1][]{boxbase,colback=excol!5,colframe=excol,
  title={\textsc{Exemple}\ifx\relax#1\relax\relax\else\textmd{ : \itshape #1}\fi}}
\newtcolorbox{piege}[1][]{boxbase,colback=attncol!6,colframe=attncol,
  title={\textsc{Piège}\ifx\relax#1\relax\relax\else\textmd{ : \itshape #1}\fi}}
\newtcolorbox{remarque}[1][]{boxbase,colback=black!4,colframe=black!45,
  title={\textsc{Remarque}\ifx\relax#1\relax\relax\else\textmd{ : \itshape #1}\fi}}
% --- boîtes EXERCICE / CORRIGÉ (noms EXACTS, auto-comptées) ---
\newtcolorbox[auto counter]{exo}[1][]{boxbase,colback=primary!4,colframe=primary,
  title={\textsc{Exercice}~\thetcbcounter\ifx\relax#1\relax\relax\else\textmd{ --- \itshape #1}\fi}}
\newtcolorbox[auto counter]{corrige}[1][]{boxbase,colback=excol!4,colframe=excol,
  title={\textsc{Corrigé}~\thetcbcounter\ifx\relax#1\relax\relax\else\textmd{ --- \itshape #1}\fi}}
\newcommand{\R}{\mathbb{R}}\newcommand{\N}{\mathbb{N}}\newcommand{\Z}{\mathbb{Z}}\newcommand{\ds}{\displaystyle}
% (un \usepackage{fancyhdr}+en-tête est possible ; il est ignoré côté web)
% =========================================================================
\begin{document}
\sisetup{per-mode=symbol}
\begin{exo}[l'ambiguïté des entiers]
Ces mesures sont écrites sous forme d'un entier, sans virgule : leur nombre de chiffres significatifs
est \emph{ambigu}. Pour chacune, indique les nombres de CS possibles, puis réécris-la en
\textbf{notation scientifique} pour qu'elle possède exactement la précision demandée.
\begin{enumerate}[label=\alph*)]
  \item $200$ \quad (à écrire avec $3$ CS)
  \item $4500$ \quad (à écrire avec $2$ CS)
  \item $70$ \quad (à écrire avec $2$ CS)
\end{enumerate}
\end{exo}

\begin{exo}[lire les CS sur la notation scientifique]
En notation scientifique, \emph{tout ce qui est écrit dans la mantisse} est significatif. Donne le
nombre de chiffres significatifs.
\begin{enumerate}[label=\alph*)]
  \item $4,00\times10^{3}$
  \item $1,0\times10^{-7}$
  \item $6,022\times10^{23}$
  \item $2,50\times10^{-2}$
\end{enumerate}
\end{exo}

\begin{exo}[le changement d'unité ne change pas les CS]
Une balance affiche la masse $m=\qty{13,02}{\gram}$. On réécrit \emph{cette même masse} dans d'autres
unités. Pour chaque écriture, donne le nombre de CS et conclus.
\begin{enumerate}[label=\alph*)]
  \item $\num{1,302e4}\ \si{\milli\gram}$
  \item $\qty{0,01302}{\kilo\gram}$
  \item $\num{1,302e-5}\ \si{\tonne}$
\end{enumerate}
\end{exo}

\begin{exo}[mesuré ou exact ?]
Pour chaque nombre, dis s'il est \textbf{mesuré} (nombre \emph{fini} de CS, à préciser) ou
\textbf{exact} (\emph{infinité} de CS).
\begin{enumerate}[label=\alph*)]
  \item les $2$ atomes d'hydrogène d'une molécule d'\ce{H2O}
  \item un volume lu $\qty{24,15}{\milli\litre}$ sur une burette
  \item la relation $\qty{1}{\hour}=\qty{3600}{\second}$
  \item une masse pesée $\qty{0,250}{\gram}$
\end{enumerate}
\end{exo}

\begin{exo}[comparer la précision]
Voici trois mesures d'une \emph{même} longueur (environ $\qty{45}{\milli\metre}$), exprimées dans des
unités différentes. Range-les de la \textbf{moins} précise à la \textbf{plus} précise, en justifiant
par le nombre de CS.
\begin{enumerate}[label=\alph*)]
  \item $\qty{0,04500}{\metre}$
  \item $\qty{4,5}{\centi\metre}$
  \item $\qty{45,0}{\milli\metre}$
\end{enumerate}
\end{exo}

\end{document}
